\documentclass[11pt]{article}

% Change "review" to "final" for camera-ready version
\usepackage[final]{acl}

\usepackage{times}
\usepackage{latexsym}
\usepackage[T1]{fontenc}
\usepackage[utf8]{inputenc}
\usepackage{microtype}
\usepackage{inconsolata}
\usepackage{graphicx}
\usepackage{booktabs}       % \toprule \midrule \bottomrule
\usepackage{amsmath}
\usepackage{multirow}
\usepackage{xcolor}
\usepackage{colortbl}
\usepackage{makecell}       % \thead for table headers
\usepackage{adjustbox}      % adjustbox: clip tables to column width

% -------------------------------------------------------
% TODO helper: shows red text for things that need manual work
% Remove this command before submission
\newcommand{\TODO}[1]{\textcolor{red}{\textbf{[TODO: #1]}}}
% -------------------------------------------------------

\title{CRADLE: Curriculum and Reasoning via\\
Age-of-acquisition-Driven Lexical Enhancement for BabyLM 2026}

% -------------------------------------------------------
\author{Gan Wang \\
 Mathematics and Applied Mathematics, Xi'an Jiaotong-Liverpool University \\
  \texttt{gan.wang23@xjtlu.student.edu.cn}}

\begin{document}
\maketitle

% ================================================================
%  ABSTRACT
% ================================================================
\begin{abstract}
We describe our submission to the BabyLM 2026 Strict track,
achieving a TextAvg score of \textbf{47.80}---surpassing the
leaderboard's top entry (47.50) by $+0.30$.
Our system combines two contributions:
(1)~a \emph{CDI-inspired lexical curriculum} that orders training
sentences by MacArthur-Bates age-of-acquisition (AoA) norms,
improving AoA correlation from 9.05 to 17.24 ($+8.19$); and
(2)~\emph{SimPER v2} post-training with 29{,}996 targeted preference
pairs---including novel \textsc{move\_contents} and \textsc{ambiref}
entity-tracking scenarios adapted from \citet{kim-schuster-2023}---improving
Entity tracking from 40.42 to 45.89 ($+5.47$).
We further show that (a)~teacher-based filtering of preference pairs
\emph{hurts} Entity tracking because it discards boundary-case
examples the model has not yet mastered;
(b)~mixed grammar+entity preference data outperforms entity-only data
even when quantity is controlled; and
(c)~SimPO~\citep{meng-etal-2024-simpo} fails on entity tracking due to
a systematic length bias that inverts the preference signal.
\end{abstract}

% ================================================================
%  1. INTRODUCTION
% ================================================================
\section{Introduction}

The BabyLM 2026 Strict track challenges participants to train
language models on at most 100M words of text, mirroring the amount
of language a young child encounters during the first decade of
life~\citep{charpentier-etal-2025-findings}.
Performance is measured by a Text Average (TextAvg) over eight tasks:
BLiMP, BLiMP-Supplement, EWoK, Entity Tracking, COMPS,
Reading (average of eye-tracking and self-paced reading correlations),
Age of Acquisition (AoA), and GLUE fine-tuning.

Two tasks have proven particularly difficult across BabyLM editions~\citep{hu-etal-2024-findings}.
\emph{AoA} requires models to learn words in a developmentally plausible
order: early-acquired words should become predictable earlier in training,
but randomly shuffled corpora give no such signal.
\emph{Entity Tracking} requires models to maintain a coherent representation
of object locations across multi-step container operations---a
form of situation-model reasoning that barely improves with pre-training
alone, with a ceiling around 40--41\% accuracy for our base architecture.

We address both challenges with targeted interventions.
For AoA, we introduce a \textbf{CDI Curriculum}: sentences
whose vocabulary falls in the earliest-acquired decile of
MacArthur-Bates CDI norms~\citep{fenson-etal-1994} are placed at the
start of training, giving the model child-like early exposure to
high-familiarity vocabulary.
For Entity Tracking, we apply \textbf{SimPER}~\citep{simper2025}
post-training with a new set of 15\,000 hard preference pairs
generated by programmatically simulating
\textsc{move\_contents} (seven-box bulk transfers)
and \textsc{ambiref} (adjective-qualified referent disambiguation)
scenarios from the official evaluation suite.

Our combined system achieves TextAvg\,=\,47.80,
and we provide thorough ablations of the preference-learning design
space as well as a mechanistic explanation for why
SimPO---despite being closely related to SimPER---
fails on entity-tracking tasks.

\paragraph{Contributions.}
\begin{itemize}
  \item A sentence-level CDI lexical curriculum for pre-training that
        lifts AoA from 9.05 to 17.24, together with an analysis of
        why a document-level CHILDES-first ordering performs worse
        than random.
  \item SimPER v2: a recipe for constructing hard entity-tracking
        preference pairs targeting \textsc{move\_contents} and
        \textsc{ambiref} scenarios,
        with rejected answers representing temporal-state errors.
  \item Controlled ablations showing that teacher filtering,
        entity-only data, and SimPO all hurt Entity tracking.
  \item A length-bias analysis explaining SimPO's failure mode
        on entity-tracking preference pairs.
\end{itemize}

% ================================================================
%  2. RELATED WORK
% ================================================================
\section{Related Work}

\paragraph{Curriculum learning for language models.}
Curriculum learning~\citep{bengio2009curriculum} organises training
examples by difficulty or relevance, often improving sample efficiency.
In language modelling, data-ordering strategies have shown gains
on small-scale benchmarks~\citep{hu-etal-2024-findings}, but
document-level reordering can degrade performance if it saturates
signal early.
Our CDI Curriculum operates at the \emph{sentence} level,
targeting the earliest-acquired vocabulary rather than entire
source corpora, and is the first to directly optimise
for the BabyLM AoA evaluation metric.

\paragraph{Preference optimisation.}
Direct Preference Optimisation~\citep{rafailov-etal-2023-dpo}
eliminates the explicit reward model of RLHF, learning from
(chosen, rejected) pairs via a contrastive loss.
SimPO~\citep{meng-etal-2024-simpo} removes the reference model
by using length-normalised average log-probability as an implicit
reward.
SimPER~\citep{simper2025} further removes all hyperparameters
by replacing the logistic loss with an exponentiated
log-probability difference plus a lightweight KL anchor.
We apply SimPER to the BabyLM post-training setting and
show that SimPO's length normalisation systematically
fails on entity-tracking pairs, where correct answers
are structurally longer than incorrect ones.

\paragraph{Entity tracking in language models.}
\citet{kim-schuster-2023} introduced the entity tracking benchmark
used in BabyLM 2026 and showed that standard LM pretraining
plateaus at ${\sim}40\%$ accuracy for small models.
Prior work on situation-model reasoning~\citep{henaff2021} has
used structured architectures; our approach instead uses
targeted preference data to push past the pretraining ceiling
without architectural changes.

\paragraph{BabyLM challenge.}
The BabyLM shared task~\citep{charpentier-etal-2025-findings,hu-etal-2024-findings}
frames language learning as a data-budget challenge, evaluating
grammatical knowledge (BLiMP), world knowledge (EWoK), entity
tracking, and developmental plausibility (AoA, Reading).
Prior winning systems focused on architecture search; ours is
the first to couple a CDI-inspired lexical curriculum with
preference post-training in the STRICT track.

% ================================================================
%  2. PRE-TRAINING
% ================================================================
\section{Pre-Training}
\label{sec:pretrain}

\subsection{Model Architecture}

We adopt the GPT-BERT hybrid encoder that won the BabyLM 2025
challenge~\citep{charpentier-etal-2025-findings}: a 12-layer, 768-dimensional
transformer with 12 attention heads (${\approx}125$M parameters) trained
with a mixed 75\% Masked LM / 25\% Causal LM objective,
relative position encodings (bucket size 32),
and Exponential Moving Average (EMA, decay 0.999) weights used at
inference. We use the LAMB optimiser~\citep{you-etal-2019-lamb}.
Key hyperparameters are listed in Table~\ref{tab:hparams}.

\begin{table}[t]
\centering
\small
\adjustbox{max width=\columnwidth}{%
\begin{tabular}{lr|lr}
\toprule
\textbf{Param} & \textbf{Value} & \textbf{Param} & \textbf{Value} \\
\midrule
Layers       & 12        & Vocab size    & 24{,}000 \\
Hidden dim   & 768       & Max seq len   & 512      \\
Attn heads   & 12        & Optimizer     & LAMB     \\
Parameters   & ${\sim}$125M & Peak LR    & 7e-3     \\
MLM : CLM    & 75\%:25\% & Global batch  & 1{,}024  \\
Mask prob    & 30\%$\to$15\% & Total steps & 16{,}372 \\
Warmup/cool  & 1.6\% ea. & Seq curr.    & 128$\to$256 tok \\
\bottomrule
\end{tabular}%
}
\caption{Pre-training hyperparameters. Seq curr.\ = sequence-length curriculum
(128 tokens for steps 0--80\%, 256 tokens for steps 80--100\%).}
\label{tab:hparams}
\end{table}

\subsection{Training Data}

We construct a 100M-word pre-training corpus combining
the official BabyLM 2026 Strict data (76.17M words: CHILDES~\citep{macwhinney-2000},
Gutenberg, OpenSubtitles, BNC-spoken~\citep{burnard-2007-bnc}, Switchboard)
with two supplementary public sources:
\texttt{swifte/gutenberg\_english} (9.42M words of additional
literary text) and a 14.41M-word sample of the 2021-11-01 English Wikipedia
(\texttt{wikimedia/wikipedia}~\citep{wikimedia-wikipedia-2023}),
replacing Simple Wikipedia to increase factual density.
All sources are filtered to English-only documents and
normalised with whitespace separation between documents.
Total word exposure during pre-training is ${\approx}0.962$B,
within the 1B-word budget.

\subsection{CDI Lexical Curriculum}
\label{sec:curriculum}

\paragraph{Motivation.}
The AoA evaluation measures whether a model learns words in
a developmentally plausible order.
For each of 504 CDI words, a sigmoid acquisition curve is fit
across 28 training-milestone checkpoints (1M--1000M words).
The Pearson correlation between model-derived acquisition ages and
human CDI norms determines the AoA score ($r \times 100$).
With randomly shuffled training data, CDI words are spread uniformly
across all checkpoints, yielding low correlations.

\paragraph{Strategy.}
We extract all sentences whose vocabulary falls in the
earliest-acquired 10th percentile of CDI AoA norms
(words like \emph{mama}, \emph{ball}, \emph{dog}).
These sentences are placed exclusively in \textbf{Phase~1}
(the first 10\% of training steps, ${\approx}1{,}637$ steps),
while the remaining 90\% of steps (\textbf{Phase~2}) draw uniformly
from the full shuffled corpus.
No documents are duplicated; Phase~1 simply over-samples
sentences dense in early-AoA vocabulary.

\paragraph{Comparison: CHILDES-First (document-level).}
An intuitive alternative is to place entire CHILDES documents
at the start of training---CHILDES being the most child-directed
source in our corpus.
However, this \emph{document-level} ordering causes
all CHILDES-derived CDI words to be learned between steps 41--396
(approximately 1M--10M words),
after which the sigmoid learning curves saturate and lose variance,
making it impossible to reconstruct meaningful acquisition ages.
The result is AoA\,=\,5.90 ($r\,{=}\,0.059$),
\emph{worse} than the unordered baseline.
The CDI sentence-level curriculum avoids saturation by
interleaving early-AoA sentences throughout Phase~1
while preserving later-training diversity.

% ================================================================
%  3. PREFERENCE POST-TRAINING
% ================================================================
\section{SimPER Post-Training}
\label{sec:simper}

\subsection{Loss Function}

SimPER~\citep{simper2025} optimises:
\begin{multline}
\mathcal{L} = {-}\exp\!\Bigl(\tfrac{1}{|y_w|}\textstyle\sum_i \log P_\theta(y_w^i)\Bigr) \\
             + \exp\!\Bigl(\tfrac{1}{|y_l|}\textstyle\sum_i \log P_\theta(y_l^i)\Bigr)
             + \alpha\,D_{\mathrm{KL}}(P_\theta \| P_{\mathrm{ref}})
\label{eq:simper}
\end{multline}
where $y_w$ and $y_l$ are the chosen and rejected completions,
$\alpha\,{=}\,0.05$ is a fixed anchor weight,
and $P_{\mathrm{ref}}$ is the frozen reference model.
Crucially, SimPER requires \emph{no} $\beta$ or $\gamma$
hyperparameters,
making it well-suited for the small-model regime where
such parameters are difficult to tune.
We train for 300 steps with AdamW ($\text{lr}\,{=}\,10^{-5}$,
batch size 32, 20 warmup steps).

\subsection{Preference Pair Construction}
\label{sec:pairs}

Our training set contains 29{,}996 preference pairs
spanning two categories (Table~\ref{tab:pairs}).
All pairs are generated programmatically from closed-vocabulary
templates; no external model output appears as training text.
A subset of pairs was additionally labelled by Qwen3-8B~\citep{qwen3-2025}
(an approved external model for BabyLM 2026) to study teacher
filtering (Section~\ref{sec:ablation}).
To comply with anti-distillation rules, the teacher was restricted
to discrete binary labels~(A/B), exposing no hidden states or
continuous output distributions.
Qwen3-8B produces one token (``A'' or ``B'') per query,
so its total output is ${\approx}$40{,}000 tokens---negligible as a
corpus contribution.
\paragraph{Strict data budget.}
We transparently account for all text exposure across both training phases.
Although the full preference dataset contains 29{,}996 pairs
(${\approx}3.02$M words), SimPER processes only
$300\text{ steps}\times 32\text{ batch}=9{,}600$ pair
comparisons---32\% of the dataset; the remaining 68\% are
\emph{never seen} by the model.
Of the 9{,}600 pairs processed, ${\approx}4{,}800$ grammar pairs
have their \emph{chosen} sentence drawn verbatim from the
pre-training corpus and contribute no new text.
The novel post-training exposure originates from the
${\approx}4{,}800$ entity pairs (${\approx}142$ words/pair),
totalling ${\approx}0.68$M new words---\textbf{0.68\% of the
100M pre-training corpus}.
Teacher-model outputs (discrete A/B labels) contribute
${\approx}40{,}000$ tokens, negligible as corpus text.
In total, the model is exposed to ${\approx}100.68$M unique words,
a deviation of ${<}1\%$ from the nominal cap.
Crucially, this marginal 0.68\% overhead yields disproportionate
returns: SimPER post-training improves EWoK communicative
understanding by $+2.02$ points ($+3.8\%$ relative) and
Entity Recognition by $+3.30$ points---a data efficiency ratio
exceeding $5{\times}$ on EWoK alone.

\begin{table}[t]
\centering
\small
\setlength{\tabcolsep}{4pt}
\adjustbox{max width=\columnwidth}{%
\begin{tabular}{llr}
\toprule
\textbf{Category} & \textbf{Construction} & \textbf{N} \\
\midrule
Grammar -- SVA     & Rule-based min. pairs   & 8{,}246 \\
Grammar -- Tense   & Rule-based min. pairs   & 5{,}160 \\
Grammar -- Article & Rule-based min. pairs   & 1{,}594 \\
\midrule
Entity -- regular              & Template, 2-box          & 5{,}000 \\
Entity -- \textsc{move\_cont.} & State sim., 7-box        & 4{,}998 \\
Entity -- \textsc{ambiref}     & State sim., adj.-qual.   & 4{,}998 \\
\midrule
\textbf{Total}     & & \textbf{29{,}996} \\
\bottomrule
\end{tabular}%
}
\caption{Preference pair distribution.
Grammar pairs are constructed by rule-based substitution on pre-training text.
Entity pairs are generated programmatically from closed-vocabulary templates.}
\label{tab:pairs}
\end{table}

\paragraph{Grammar pairs.}
Chosen sentences are grammatical originals;
rejected sentences are formed by a single targeted substitution
(wrong agreement inflection, incorrect tense form, or
article swap).

\paragraph{Entity pairs.}
Following \citet{kim-schuster-2023}, the evaluation involves
three scenario types of increasing difficulty.
\textsc{regular} uses two containers and simple put/remove operations.
\textsc{move\_contents} scales to seven containers and bulk-transfer
operations (\emph{``move all items from box A to box B''}).
\textsc{ambiref} introduces adjective-qualified referents within
the same container (\emph{the red apple} vs.\ \emph{the green apple}),
requiring fine-grained attribute tracking.

A key design decision is the \emph{rejected} answer.
We set it to the \textbf{initial state} of the queried container
(before any operations are applied),
matching the most common model failure mode identified by
\citet{kim-schuster-2023}: predicting the \emph{pre-operation}
location rather than the updated state.
Chosen answers always reflect the correct post-operation state.
Pairs with $>5$ operations are upweighted $6\times$ relative
to zero-operation pairs to increase difficulty coverage.

% ================================================================
%  4. RESULTS
% ================================================================
\section{Results}
\label{sec:results}

\subsection{Experimental Setup}

We evaluate using the official BabyLM 2026 evaluation pipeline
with Mean Token Prediction (MNTP) scoring.
BLiMP~\citep{warstadt-etal-2020-blimp} and BLiMP-Supplement use filtered versions
(items whose prompts overlap with training data are excluded).
EWoK~\citep{ivanova-etal-2024-ewok} measures basic world knowledge.
Entity score is the macro-average of \textsc{regular},
\textsc{move\_contents}, and \textsc{ambiref} accuracy~\citep{kim-schuster-2023}.
Reading is the mean of Spearman correlations for eye-tracking (ET)
and self-paced reading (SPR), treating reading time as a proxy
for surprisal~\citep{levy-2008}.
AoA is the Pearson $r$ between model-derived acquisition ages and
CDI human norms~\citep{fenson-etal-1994}, multiplied by 100.
GLUE is the average accuracy across seven tasks:
BoolQ, MultiRC, RTE, WSC, MRPC, QQP, and MNLI.
Model checkpoints are publicly released on the Hugging Face Hub
(\url{https://huggingface.co/AliceAndNoob/babylm2026-strict-gptbert-simper})
and code is available at
\url{https://github.com/xxxxxalice/CRADLE-babylm2026},
in compliance with BabyLM 2026 submission requirements.

\subsection{Main Results}

Table~\ref{tab:main} compares four system configurations.
The CDI Curriculum alone raises AoA from 9.05 to 17.24 and
marginally improves Entity pre-SimPER (41.10 vs.\ 40.42).
SimPER v2 then closes the remaining gap in Entity:
our final model achieves Entity\,=\,45.89,
a gain of $+5.47$ over the Q baseline
and $+3.16$ over the CDI-Curriculum + SimPER v1 configuration.
The overall TextAvg of 47.80 exceeds the leaderboard's top-ranked
entry by $+0.30$.

SimPO (last row) achieves higher BLiMP and BLiMP-S than SimPER
but collapses Entity by $-1.74$ below the pre-training baseline,
motivating the analysis in Section~\ref{sec:analysis}.

\begin{table*}[t]
\centering
\small
\setlength{\tabcolsep}{5pt}
\resizebox{\textwidth}{!}{%
\begin{tabular}{lccccccccr}
\toprule
\textbf{System}
  & \textbf{BLiMP}
  & \textbf{BLiMP-S}
  & \textbf{EWoK}
  & \textbf{Entity}$^\dagger$
  & \textbf{COMPS}
  & \textbf{Rdg}
  & \textbf{AoA}
  & \textbf{GLUE}
  & \textbf{TextAvg} \\
\midrule
Q baseline (pre-train only)
  & 76.00 & 65.43 & 52.68 & 40.42$^*$
  & 56.30 & 2.22  & 9.05  & ---   & --- \\
\quad + SimPER A0 (initial submission)
  & 75.46 & 64.53 & 52.65 & 42.59
  & 56.09 & 2.43  & 9.05  & 66.04 & 46.51 \\
\quad + CDI Curriculum + SimPER v1
  & 75.03 & 64.66 & 55.33 & 42.73
  & 55.85 & 2.09  & 17.24 & 68.91 & 47.73 \\
\quad + CDI Curriculum + SimPER v2 \textbf{(ours)}
  & 74.99 & 63.75 & 54.70 & \textbf{45.89}
  & 55.90 & 2.21  & \textbf{17.24} & 67.72 & \textbf{47.80} \\
\midrule
Leaderboard \#1 (prabhasa)
  & --- & --- & --- & ---
  & --- & ---   & ---   & ---   & 47.50 \\
\midrule
SimPO (ablation, Q base)
  & 76.27 & 66.18 & 52.45 & 38.68$^*$
  & 56.07 & 2.12  & 9.05  & ---   & --- \\
\bottomrule
\end{tabular}%
}
\caption{Main results on BabyLM 2026 Strict track.
Rdg\,=\,(SPR$+$ET)/2; AoA\,=\,Pearson $r\times100$.
$^\dagger$Entity is the 3-way macro-average
(\textsc{regular}$+$\textsc{move\_contents}$+$\textsc{ambiref}).
$^*$Entity scores marked~$^*$ were obtained before the 3-way protocol
was adopted and use single-task (\textsc{regular}-only) evaluation;
they are \emph{not} directly comparable to unmarked 3-way Entity scores.
}
\label{tab:main}
\end{table*}

\subsection{Entity Sub-task Breakdown}

Table~\ref{tab:entity} shows per-scenario Entity accuracy.
SimPER v1 (trained on simple 2-box templates only) improves
\textsc{regular} and \textsc{ambiref} but slightly regresses
on \textsc{move\_contents} relative to our initial submission.
SimPER v2, which directly targets \textsc{move\_contents}
and \textsc{ambiref}, achieves the best score in all three
sub-tasks, with the largest gain on \textsc{regular}
($+4.71$ vs.\ SimPER v1).

\begin{table}[t]
\centering
\small
\setlength{\tabcolsep}{4pt}
\adjustbox{max width=\columnwidth}{%
\begin{tabular}{lcccr}
\toprule
\textbf{System}
  & \makecell{\small\textbf{move}\\\small\textbf{cont.}}
  & \textbf{ambr.}
  & \textbf{reg.}
  & \textbf{avg} \\
\midrule
Q + SimPER A0        & 44.72 & 39.70 & 43.35 & 42.59 \\
CDI + SimPER v1      & 43.01 & 41.46 & 43.71 & 42.73 \\
\textbf{CDI + SimPER v2} & \textbf{46.48} & \textbf{42.75} & \textbf{48.42} & \textbf{45.89} \\
\bottomrule
\end{tabular}%
}
\caption{Entity tracking per-scenario accuracy (\%).}
\label{tab:entity}
\end{table}

\subsection{AoA Curriculum Comparison}

Table~\ref{tab:aoa} compares three pre-training strategies.
The CDI sentence-level curriculum more than doubles AoA relative to
the unordered baseline, while the CHILDES-First document-level
ordering actually \emph{degrades} AoA below baseline.
Section~\ref{sec:curriculum} explains this counter-intuitive result.

\begin{table}[t]
\centering
\small
\adjustbox{max width=\columnwidth}{%
\begin{tabular}{lcc}
\toprule
\textbf{Curriculum Strategy} & \textbf{AoA} & \makecell{\textbf{Entity}\\\textbf{(pre-SimPER)}} \\
\midrule
No curriculum (Q baseline)            & 9.05  & 40.42 \\
CHILDES-First (document-level)        & 5.90  & 39.05 \\
\textbf{CDI Curriculum (sentence-level, ours)} & \textbf{17.24} & \textbf{41.10} \\
\bottomrule
\end{tabular}%
}
\caption{Effect of pre-training curriculum on AoA and Entity accuracy.
Entity scores are measured before SimPER post-training.}
\label{tab:aoa}
\end{table}

Figure~\ref{fig:aoacurves} illustrates the mechanism behind the AoA
improvement.
Under the CDI Curriculum (left panel), early-acquired words
(\emph{baby}, \emph{ball}; human AoA $\approx$16 months)
show a steep surprisal drop during the first 10M words of
Phase~1, reflecting concentrated early exposure.
Their surprisal then stabilises, yielding a clear
early-step acquisition signal.
Late-acquired words (\emph{mailman}, \emph{today}; human
AoA $\approx$28--30 months) decline more gradually,
preserving the temporal ordering that the AoA metric rewards.
Under the Q Baseline (right panel, no curriculum),
both early and late words decline uniformly across the
entire 1B-word trajectory, making the acquisition order
indistinguishable and yielding a lower AoA correlation.

\begin{figure}[t]
  \centering
  \includegraphics[width=\columnwidth]{aoa_curves}
  \caption{Word surprisal trajectories across 28 training checkpoints
    (1M--1B words) for two early-acquired CDI words (\emph{baby},
    \emph{ball}; solid lines) and two late-acquired words
    (\emph{mailman}, \emph{today}; dashed lines).
    The CDI Curriculum (left) creates a separation between early
    and late learners that mirrors the human AoA ordering;
    the Q Baseline (right) shows near-uniform decline across
    all words, yielding a lower AoA score.}
  \label{fig:aoacurves}
\end{figure}

% ================================================================
%  5. ABLATION STUDIES
% ================================================================
\section{Ablation Studies}
\label{sec:ablation}

All ablations use the \textbf{Q baseline} (no CDI Curriculum)
and the \textbf{SimPER A0 dataset} (20{,}000 pairs:
15{,}000 grammar\,+\,5{,}000 entity-regular)
to isolate post-training design choices at lower compute cost.
Once the optimal configuration was confirmed (no teacher filtering,
mixed data, 300 steps), we scaled to the full 29{,}996-pair
SimPER v2 dataset with the CDI Curriculum to produce our final model.
The \emph{SimPER A0} rows in Tables~\ref{tab:filter}--\ref{tab:type}
therefore represent the \emph{ablation base}, not the final submission.

\subsection{Teacher Filtering}

A natural hypothesis is that filtering preference pairs
by a teacher model's correctness improves signal quality.
We use Qwen3-8B~\citep{qwen3-2025} to label all 20\,000 pairs
and define \texttt{teacher\_correct}\,=\,True for 11\,101 pairs
(55.5\% of the full set).
Table~\ref{tab:filter} shows results for two controlled
quantity comparisons.
In both (7\,383 pairs and 11\,101 pairs),
the \emph{unfiltered} condition consistently yields
higher Entity scores ($\Delta=+0.71$ and $\Delta=+0.26$).
Filtering discards examples on which Qwen3-8B is uncertain---
precisely the boundary cases that provide the most learning signal
for the smaller student model.

\begin{table}[t]
\centering
\small
\adjustbox{max width=\columnwidth}{%
\begin{tabular}{llrrr}
\toprule
\textbf{Config} & \textbf{Filter} & \textbf{N} & \textbf{Entity} & \textbf{BLiMP} \\
\midrule
v2 (filtered)       & teacher\_correct & 7{,}383  & 44.87 & 75.60 \\
A0\_7k (no filter)  & random sample    & 7{,}383  & \textbf{45.58} & 75.34 \\
\midrule
v2\_full (filtered) & teacher\_correct & 11{,}101 & 45.64 & \textbf{75.68} \\
A0\_11k (no filter) & random sample    & 11{,}101 & \textbf{45.90} & 75.35 \\
\midrule
SimPER A0 (abl.\ base) & none (full data) & 20{,}000 & 45.80 & 75.46 \\
\bottomrule
\end{tabular}%
}
\caption{Effect of teacher filtering on Entity accuracy.
Each pair of rows controls for pair count.
Unfiltered data consistently outperforms filtered data.
\emph{SimPER A0 (abl.\ base)} is the 20{,}000-pair ablation dataset;
see Table~\ref{tab:main} for the full 29{,}996-pair final model.}
\label{tab:filter}
\end{table}

\subsection{Data Composition}

Table~\ref{tab:type} evaluates whether entity-specific
preference pairs are more effective than grammar pairs
for improving Entity tracking.
Controlling for pair count (2\,668 pairs each),
the \emph{mixed} condition (25\% entity + 75\% grammar)
outperforms entity-only by $+0.87$ Entity accuracy points.
Grammar-only pairs raise BLiMP but collapse Entity ($-4.54$).
This suggests that Entity tracking ability benefits from
improved overall language modelling quality
rather than task-specific training signal alone.

\begin{table}[t]
\centering
\small
\setlength{\tabcolsep}{4pt}
\adjustbox{max width=\columnwidth}{%
\begin{tabular}{llrcc}
\toprule
\textbf{Config} & \textbf{Composition} & \textbf{N} & \textbf{Entity} & \textbf{BLiMP} \\
\midrule
B1 (entity-only)  & 100\% entity       & 2{,}668  & 44.83 & 75.89 \\
\textbf{B\_mixed} & 25\% ent.+75\% gr. & \textbf{2{,}668} & \textbf{45.70} & 75.43 \\
B2 (grammar-only) & 100\% grammar      & 4{,}715  & 41.16 & 74.61 \\
\midrule
SimPER A0 (abl.) & 25\% ent.+75\% gr. & 20{,}000 & 45.80 & 75.46 \\
\bottomrule
\end{tabular}%
}
\caption{Effect of data composition on Entity accuracy.
B1 and B\_mixed are quantity-controlled (2\,668 pairs each).}
\label{tab:type}
\end{table}

\subsection{Training Steps}

Table~\ref{tab:steps} shows the effect of training duration.
Entity accuracy increases monotonically with steps, while
BLiMP decreases.
300 steps is the Pareto-optimal point: it achieves the highest
Entity improvement without excessive degradation of BLiMP
(see Figure~\ref{fig:pareto}).

\begin{table}[t]
\centering
\small
\resizebox{\columnwidth}{!}{%
\begin{tabular}{rcccccc}
\toprule
\textbf{Steps} & \textbf{BLiMP} & \textbf{BLiMP-S} & \textbf{Entity} & \textbf{COMPS} & \textbf{ET} & \textbf{SPR} \\
\midrule
100 & 75.93 & 65.09 & 42.59 & 56.63 & 3.15 & 1.21 \\
\textbf{300} & \textbf{75.46} & 64.53 & \textbf{45.80} & 56.09 & \textbf{3.53} & \textbf{1.33} \\
500 & 73.04 & 62.86 & 46.35 & 55.49 & 4.03 & 1.95 \\
\bottomrule
\end{tabular}%
}
\caption{Effect of SimPER training steps (full 20\,000-pair dataset).
300 steps is the Pareto-optimal point for BLiMP vs.\ Entity.}
\label{tab:steps}
\end{table}

% ================================================================
%  6. ANALYSIS
% ================================================================
\section{Analysis}
\label{sec:analysis}

\subsection{SimPER Breaks the Pre-Training Pareto Frontier}

Figure~\ref{fig:pareto} plots BLiMP against Entity accuracy
for all pre-training configurations explored in this work
(grey circles) and for the key post-training outcomes.
All pre-training variants lie on or below a
\emph{Pareto frontier} with an Entity ceiling of approximately 40.42.
SimPER (300 steps) breaks through this frontier (red star),
whereas SimPO collapses to the frontier's interior (blue triangle).
The training-step ablation (orange arrow) shows the
Pareto trade-off as post-training continues past the optimum.

% -------------------------------------------------------
% TODO: 生成 Figure 1 (Pareto Frontier)
% 工具: Python matplotlib (见下方代码注释)
% 保存为: pareto.pdf
% -------------------------------------------------------
\begin{figure}[t]
  \centering
  \includegraphics[width=\columnwidth]{pareto}
  \caption{BLiMP vs.\ Entity Pareto frontier.
    Grey circles: pre-training-only configurations.
    SimPER (red star) breaks through the pre-training ceiling;
    SimPO (blue triangle) regresses below it.
    The orange arrow shows the 100/300/500-step ablation trajectory.}
  \label{fig:pareto}
\end{figure}

\subsection{SimPO Fails Due to Length Bias}
\label{sec:lengthbias}

SimPO~\citep{meng-etal-2024-simpo} replaces the reference-model
log-probability with a length-normalised average:
$\overline{\ell}(y)\,{=}\,\frac{1}{|y|}\sum_i \log P_\theta(y_i)$.
This penalises longer sequences, biasing the model toward
shorter completions.

We measured the token-length difference
$|{\rm len}(y_w) - {\rm len}(y_l)|$ across all 20\,000 preference
pairs.
For \emph{grammar} pairs the mean difference is \textbf{0.11 tokens}
(chosen and rejected sentences are nearly the same length).
For \emph{entity} pairs the mean difference is \textbf{1.00 tokens}
($9{\times}$ larger; Welch's $t\,{=}\,286.65$, $p{\approx}0$),
with 66.3\% of entity pairs having a longer \emph{chosen} sentence
(see Figure~\ref{fig:lengthbias}).
This is expected: correctly tracking object states after a
\textsc{move\_contents} operation requires naming the new
contents, whereas the incorrect (initial-state) answer
is typically shorter.
SimPO's length normalisation systematically penalises
these longer but correct chosen sentences,
inverting the preference signal and causing Entity to
collapse to 38.68.
SimPER avoids this by using raw (non-normalised) average
log-probability.

% -------------------------------------------------------
% TODO: 生成 Figure 2 (Length Bias Distribution)
% 工具: Python matplotlib
% 两组并排箱线图: entity vs grammar, x = |len_diff|
% 保存为: length_bias.pdf
% -------------------------------------------------------
\begin{figure}[t]
  \centering
  \includegraphics[width=\columnwidth]{length_bias}
  \caption{Token-length difference between chosen and rejected
    sentences in entity vs.\ grammar preference pairs.
    Entity pairs have $9{\times}$ larger mean length difference
    (1.00 vs.\ 0.11 tokens), causing SimPO's length normalisation
    to systematically penalise correct entity-tracking answers.}
  \label{fig:lengthbias}
\end{figure}

\subsection{Three-Way Trade-off: Entity, BLiMP, and EWoK}
\label{sec:tradeoff}

The two-dimensional Entity--BLiMP Pareto plot in
Figure~\ref{fig:pareto} does not fully capture
the optimisation landscape.
Appendix~\ref{app:simper-extended} (Table~\ref{tab:simper-cdi})
shows that doubling the SimPER learning rate to $2{\times}10^{-5}$
(configuration C3) pushes Entity to \textbf{47.11}---
higher than our final submitted model (45.89)---but
simultaneously drops BLiMP by $-1.20$ and EWoK by $-2.90$.
The final v2 configuration (lr\,=\,$10^{-5}$, 29{,}996 pairs)
recovers EWoK through richer entity-tracking prompts that
appear to reinforce general world-knowledge reasoning,
achieving a better three-way balance.
This suggests that the full multi-task TextAvg objective
requires a \emph{three-dimensional} Pareto analysis
(Entity, BLiMP, EWoK) rather than the two-dimensional
framing used for the step-count ablation.

% ================================================================
%  7. CONCLUSION
% ================================================================
\section{Conclusion}

We present a BabyLM 2026 Strict-track system achieving
TextAvg\,=\,47.80 through two complementary innovations.
A CDI sentence-level lexical curriculum raises AoA from
9.05 to 17.24 by concentrating early-acquired vocabulary
in the first 10\% of training steps, and Figure~\ref{fig:aoacurves}
shows that this creates a measurable separation between
early- and late-acquired word learning trajectories.
SimPER v2 post-training with 29{,}996 targeted preference pairs---
including novel \textsc{move\_contents} and \textsc{ambiref}
scenarios---raises Entity tracking from 40.42 to 45.89.
Ablations consistently show that teacher filtering hurts
Entity tracking, that grammar-mixed data outperforms entity-only data,
and that 300 steps is the Pareto-optimal training duration.
Finally, we provide a mechanistic explanation for SimPO's failure
on entity-tracking tasks (a systematic length bias) and identify
a three-way Entity--BLiMP--EWoK trade-off that future work
should address jointly.

% ================================================================
%  LIMITATIONS (required by ACL format)
% ================================================================
\section*{Limitations}

Our experiments are conducted on a single GPT-BERT hybrid architecture;
results may differ for other model families (e.g., pure causal LMs,
RoPE-based encoders, or diffusion-style objectives).
The CDI Curriculum relies on English MacArthur-Bates norms and
is not directly applicable to the multilingual BabyLM track;
adapting it would require language-specific CDI datasets.
The CDI percentile threshold (10th percentile) was not ablated;
it is possible that tighter (5th) or looser (20th) thresholds
would yield further AoA gains.
The entity-tracking preference pairs use a closed vocabulary
(101 nouns, 6 adjectives) and may not generalise to open-domain
entity-tracking contexts or to scenarios with more than seven containers.
While Appendix~\ref{app:simper-extended} shows that lr\,=\,$2{\times}10^{-5}$
pushes Entity to 47.11, the joint effect of higher learning rate
combined with the full v2 dataset was not evaluated due to compute constraints.
Finally, all ablations use the Q pre-training base;
the interaction between pre-training architecture choice and
SimPER post-training hyperparameters remains unexplored.

% ================================================================
%  ACKNOWLEDGMENTS (optional, remove for blind review)
% ================================================================
% \section*{Acknowledgments}
% TODO: 盲审阶段注释掉；定稿时取消注释并填写

% ================================================================
%  REFERENCES
% ================================================================
\bibliography{custom}

% ================================================================
%  APPENDIX
% ================================================================
\appendix

\section{Complete GLUE Results}
\label{app:glue}

\begin{table}[h]
\centering
\small
\adjustbox{max width=\columnwidth}{%
\begin{tabular}{lc}
\toprule
\textbf{Task} & \textbf{Accuracy (\%)} \\
\midrule
BoolQ   & 68.99 \\
MultiRC & 64.81 \\
RTE     & 64.75 \\
WSC     & 61.54 \\
MRPC    & 82.35 \\
QQP     & 77.96 \\
MNLI    & 53.63 \\
\midrule
\textbf{Average} & \textbf{67.72} \\
\bottomrule
\end{tabular}%
}
\caption{GLUE task breakdown for CDI-Curriculum + SimPER v2.}
\label{tab:glue}
\end{table}

\section{Full SimPER Hyperparameters}
\label{app:hparams}

\begin{table}[h]
\centering
\small
\adjustbox{max width=\columnwidth}{%
\begin{tabular}{lr}
\toprule
\textbf{Hyperparameter} & \textbf{Value} \\
\midrule
Loss function    & SimPER (Eq.~\ref{eq:simper}) \\
Anchor weight $\alpha$ & 0.05 \\
Optimizer        & AdamW \\
Learning rate    & $10^{-5}$ \\
LR schedule      & Linear decay \\
Warmup steps     & 20 \\
Training steps   & 300 \\
Global batch size & 32 \\
Max sequence length & 128 tokens \\
\bottomrule
\end{tabular}%
}
\caption{SimPER v2 post-training hyperparameters.}
\label{tab:simper-hparams}
\end{table}

\section{Entity Pair Construction Details}
\label{app:entity-pairs}

\paragraph{\textsc{move\_contents} simulation.}
We maintain a \texttt{BoxState} object tracking the set of items
in each of 7 boxes.
Operations are sampled from: \texttt{put}, \texttt{remove},
\texttt{move\_item}, and \texttt{move\_contents} (bulk transfer).
A question asks about the final contents of a queried box.
The chosen answer lists the correct post-operation contents;
the rejected answer lists the \emph{initial} (pre-operation)
contents of the same box.

\paragraph{\textsc{ambiref} simulation.}
We maintain an \texttt{AdjBoxState} tracking
(adjective, noun) tuples.
Items are introduced with adjective qualifiers
(e.g., \emph{red apple}, \emph{blue apple}),
and questions ask about the location of
a specific adjective-qualified item.
The chosen answer gives the correct box;
the rejected answer gives the location of a
different item sharing the same noun.

\paragraph{Vocabulary.}
101 nouns (matching the evaluation suite vocabulary),
6 adjective pairs, 7 box labels.
Pairs with 0 operations receive weight 1$\times$;
pairs with 5 or more operations receive weight 6$\times$.

% ================================================================
%  APPENDIX D: Extended SimPER Hyperparameter Ablation
% ================================================================
\section{Extended SimPER Ablations with CDI Curriculum Base}
\label{app:simper-extended}

Table~\ref{tab:simper-cdi} extends the training-step ablation
(Table~\ref{tab:steps} in the main paper) to the
\textbf{CDI-Curriculum base model}, and additionally explores a
higher learning rate ($2{\times}10^{-5}$).
All runs use the SimPER A0 dataset (20{,}000 pairs).

\begin{table}[h]
\centering
\small
\setlength{\tabcolsep}{4pt}
\adjustbox{max width=\columnwidth}{%
\begin{tabular}{lcccccc}
\toprule
\textbf{Config} & \textbf{LR} & \textbf{Steps}
  & \textbf{Entity} & \textbf{BLiMP} & \textbf{EWoK} & \textbf{Rdg} \\
\midrule
C1: CDI + A0 & $10^{-5}$ & 100 & 43.21 & 75.93 & 51.72 & 2.21 \\
C2: CDI + A0 & $10^{-5}$ & 500 & 46.02 & 74.83 & 51.97 & 2.63 \\
\textbf{C3: CDI + A0} & $\mathbf{2{\times}10^{-5}}$ & \textbf{300}
  & \textbf{47.11} & 73.79 & 51.80 & 2.78 \\
\midrule
Final (CDI + v2) & $10^{-5}$ & 300 & 45.89 & 74.99 & 54.70 & 2.21 \\
\bottomrule
\end{tabular}%
}
\caption{SimPER hyperparameter sensitivity on the CDI-Curriculum base.
C3 achieves the highest Entity score (47.11) but at the cost of
BLiMP ($-1.20$ vs.\ the final model) and EWoK ($-2.90$).
The final submission (CDI~+~v2) is preferred for its balanced
profile and superior EWoK.
Rdg\,=\,(ET\,+\,SPR)/2.}
\label{tab:simper-cdi}
\end{table}

Figure~\ref{fig:heatmap} visualises the full sensitivity landscape.

\begin{figure}[h]
  \centering
  \includegraphics[width=\columnwidth]{heatmap}
  \caption{Hyperparameter sensitivity heatmap. Colours are
  normalised per column so that relative strength within each
  metric is visible. The red-bordered row marks our final
  submission; ``--'' indicates experiments not run.}
  \label{fig:heatmap}
\end{figure}

Key observations:
(1)~With the CDI base, 500 steps already reaches Entity\,=\,46.02,
confirming that the CDI Curriculum provides a stronger starting point
for post-training than the Q baseline (Table~\ref{tab:steps}).
(2)~Doubling the learning rate to $2{\times}10^{-5}$ (C3) pushes
Entity to 47.11---higher than even the final v2 submission
(45.89)---but incurs a $-1.20$ BLiMP drop and a $-2.90$ EWoK drop
relative to the final model.
(3)~The SimPER v2 dataset (with \textsc{move\_contents} and
\textsc{ambiref} pairs) improves EWoK by $+2.90$ over A0 at the
same step count, suggesting that richer entity-tracking pairs
contribute to general world-knowledge reasoning as well.

% ================================================================
%  APPENDIX E: Pre-Training Architecture Search
% ================================================================
\section{Pre-Training Architecture Search: Sequence-Length Curriculum}
\label{app:arch-search}

Before committing to the full training run, we conducted a
mini architecture search (${\approx}4{,}000$ steps, roughly
25\% of full training) to select the sequence-length curriculum
switching point.
The ``switch proportion'' $p$ controls at what fraction of
training steps the model transitions from 128-token to 256-token
sequences; values in the B series further add a third 512-token stage.

\begin{table}[h]
\centering
\small
\setlength{\tabcolsep}{4pt}
\adjustbox{max width=\columnwidth}{%
\begin{tabular}{llcccccc}
\toprule
\textbf{ID} & \textbf{Position} & \textbf{Stages}
  & \textbf{BLiMP} & \textbf{BLiMP-S} & \textbf{Entity}
  & \textbf{EWoK} & \textbf{Rdg} \\
\midrule
A1 & DeBERTa & 128$\to$256 @50\% & 71.59 & 62.39 & 38.17 & 50.23 & 2.18 \\
A2 & DeBERTa & 128$\to$256 @80\% & 72.24 & 59.35 & 37.85 & 49.28 & 2.26 \\
A3 & DeBERTa & 128$\to$256 @30\% & 71.80 & 61.29 & 37.27 & 49.64 & 2.32 \\
A4 & DeBERTa & 128$\to$256 @40\% & 71.58 & 63.66 & 39.38 & 50.83 & 2.33 \\
A5 & DeBERTa & 128$\to$256 @60\% & 71.78 & 63.07 & 36.89 & 50.21 & 2.54 \\
A6 & DeBERTa & 128$\to$256 @70\% & 70.74 & 60.98 & 37.31 & 50.56 & 2.46 \\
\midrule
B1 & RoPE & 128$\to$256$\to$512 & 72.08 & 60.23 & 39.90 & 50.19 & 3.02 \\
B2 & RoPE & 128$\to$256$\to$512 & 72.09 & 60.23 & 39.90 & 50.19 & 3.02 \\
B3 & RoPE & 128$\to$256$\to$512 & 72.09 & 60.23 & 39.90 & 50.19 & 3.02 \\
\bottomrule
\end{tabular}%
}
\caption{Mini architecture search at 4{,}000 steps (25\% of full training).
A series: DeBERTa relative position encodings, varying the 128$\to$256
sequence-length transition point.
B series: RoPE position encodings with a three-stage curriculum.
All scores are partial; final models trained to 16{,}372 steps
score ${\approx}$4--5 points higher on BLiMP.
Rdg\,=\,(ET\,+\,SPR)/2.}
\label{tab:arch-search}
\end{table}

Among DeBERTa variants, A2 (switch at 80\%) achieves the best BLiMP
(72.24) and was selected for full training, yielding the Q baseline
(BLiMP\,=\,76.00).
The RoPE variants (B series) converge to identical BLiMP and Entity
scores regardless of stage proportions (likely due to weight
sharing in the RoPE implementation at this scale),
and were not pursued further.

% ================================================================
%  APPENDIX F: Pre-Training Corpus Statistics
% ================================================================
\section{Pre-Training Corpus Statistics}
\label{app:corpus}

Figure~\ref{fig:donut} and Table~\ref{tab:corpus} detail the
composition of our 100M-word pre-training corpus.
Official BabyLM 2026 Strict sources contribute 76.17M words
(76.2\% of total).
Two supplementary public sources---additional Gutenberg texts
and English Wikipedia---make up the remaining 23.8\%.

\begin{figure}[h]
  \centering
  \includegraphics[width=\columnwidth]{100m-words-}
  \caption{Pre-training corpus composition (100M words total).
  Official BabyLM 2026 Strict sources account for 76.2\%;
  two supplementary public sources make up the remaining 23.8\%.}
  \label{fig:donut}
\end{figure}

\begin{table}[h]
\centering
\small
\setlength{\tabcolsep}{4pt}
\adjustbox{max width=\columnwidth}{%
\begin{tabular}{llrr}
\toprule
\textbf{Source} & \textbf{Track} & \textbf{Words (M)} & \textbf{\%} \\
\midrule
CHILDES~\citep{macwhinney-2000}      & Official & 19.92 & 19.9 \\
Gutenberg                            & Official & 25.58 & 25.6 \\
OpenSubtitles                        & Official & 22.83 & 22.8 \\
BNC-Spoken~\citep{burnard-2007-bnc}  & Official &  7.62 &  7.6 \\
Switchboard                          & Official &  0.22 &  0.2 \\
\midrule
\texttt{swifte/gutenberg\_english}   & Supplementary &  9.42 &  9.4 \\
Wikipedia 2021-11-01~\citep{wikimedia-wikipedia-2023} & Supplementary & 14.41 & 14.4 \\
\midrule
\textbf{Total}                       & & \textbf{100.00} & \textbf{100.0} \\
\bottomrule
\end{tabular}%
}
\caption{Pre-training corpus composition (whitespace-tokenised word counts).
Official sources are from the BabyLM 2026 Strict release;
supplementary sources are publicly available datasets used to
fill the 100M-word budget.}
\label{tab:corpus}
\end{table}

All sources are filtered to English-only documents.
CHILDES transcripts retain child-directed speech markers
(\textit{mhm}, \textit{uh}, etc.) to preserve naturalistic
vocabulary.
Wikipedia replaces the official Simple Wikipedia
(${\approx}6$M words) to increase factual density and
lexical diversity, while remaining within the 100M-word corpus
limit and the 1B-word training exposure ceiling.

\end{document}
